

\section{Week 15: FOL Proofs}

%\subsection{Basic rules for FOL} \cite[][Chapter 36, pp. 297-315]{Mag25}

Proofs for FOL will be governed by the same rule system as TFL, with the
addition of some new rules for dealing with singular terms, predicates and
quantifiers. Of course, these rules will take the form of our familiar
introduction and elimination rules, plus some derived rules that let us avoid
repeating frequently needed subproofs.

Given the definitions for substitution instances discussed last week, defining
elimination rules for $eforall x$ and introduction rules for $eexists x$ is
particularly straight-forward.\footnote{In our rules, we let \egamma be
metavariables for FOL variables and \ealpha be a metavariable for FOL names}.

\begin{fitchrule}[Universal Elimination ($\eforall$E)]
  \label{def:universal-elimination-rule}
  \leavevmode
  \vspace{-0.9\baselineskip}
  \begin{fitchproof}
    \have[n] {1} {$\eforall\egamma\ephi(\dots\egamma\dots\egamma\dots)$}
    \have[ ] {2} {$\ephi(\dots\ealpha\dots\ealpha\dots)$} \Ae{1}
  \end{fitchproof}
\end{fitchrule}

With existential introduction, we can even substitute the variable for any one
or more instances of the relevant name, which we will depict as
$\eexists\egamma\ephi(\dots\egamma\dots\ealpha\dots)$.

\begin{fitchrule}[Existential Introduction ($\eexists$I)]
  \label{def:existential-introduction-rule}
  \leavevmode
  \vspace{-0.9\baselineskip}
  \begin{fitchproof}
    \have[n] {1} {$\ephi(\dots\ealpha\dots\ealpha\dots)$}    
    \have[ ] {2} {$\eexists\egamma\ephi(\dots\egamma\dots\ealpha\dots)$} \Ee{1}
  \end{fitchproof}
\end{fitchrule}

The two complementary operations (\emph{existential elimination} and
\emph{universal introduction}) are more contrained in their definitions. For
universal introduction ($\eforall$I), we should only allow the rule to apply
when the name in question is \emph{arbitrary}. For the purposes of proofs, this
amounts to not being found in premises or unresolved subproofs since their use
there makes them non-arbitrary.

\begin{fitchrule}[Universal Introduction ($\eforall$I)]
  \label{def:universal-introduction-rule}
  \leavevmode
  \vspace{-0.9\baselineskip}
  \begin{fitchproof}
    \have[n] {1} {$\ephi(\dots\ealpha\dots\ealpha\dots)$}    
    \have[ ] {2} {$\eforall\egamma\ephi(\dots\egamma\dots\egamma\dots)$} \Ee{1}
  \end{fitchproof}
  But only if \ealpha does not occur in a premise or an undischarged assumption.
\end{fitchrule}

\begin{fitchrule}[Existential Elimination ($\eexists$E)]
  \label{def:existential-elimination-rule}
  \leavevmode
  \vspace{-0.9\baselineskip}
  \begin{fitchproof}
    \have[n] {1} {$\eexists\egamma\ephi(\dots\egamma\dots\egamma\dots)$}
    \open
    \hypo[i] {2} {$\ephi(\dots\ealpha\dots\ealpha\dots)$} \by{AS}{}
    \have[j] {3} {\epsi}
    \close
    \have[ ] {4} {\epsi} \Ee{1,2-3}
  \end{fitchproof}
  But \ealpha must not occur in:\\
  \begin{enumerate}
  \item any assumption discharged before line $i$
  \item $\eexists\egamma\ephi(\dots\egamma\dots\egamma\dots)$
  \item \epsi
  \end{enumerate}
\end{fitchrule}

%\subsection{Proofs with quantifiers} \cite[][Chapter 37, pp. 316-22]{Mag25}

Similarly, proofs in FOL will proceed with a strategy based on the one we
devised for FOL, with the addition of some new techniques to deal with
quantifiers. Quite tersely, we can describe these new techniques as follows.

\begin{enumerate}
\item Work forward when dealing with the existential quantifier (\eexists)
\item Work backward when dealing with universal quantifiers (\eforall)
\end{enumerate}

% TODO: complete list of techniques ??

%\subsection{Conversion of quantifiers} \cite[][Chapter 38, pp. 323-5]{Mag25}

Like conjunction and disjunction, existential and universal quantifiers are
mutually-definable with the help of the negation operator.

\begin{fitchrule}[Conversion of Quantifiers ($CQ$))]
  \label{def:quantifier-conversion-rule}
  We can relate the quantifiers of first-order logic, with the following two
  pairs of complementary rules. One pair for handling a negated existential
  quantifier:\\
  \leavevmode
  %\vspace{-0.9\baselineskip}
  \twoitems{
    \begin{fitchproof}
      \have[n] {1} {$\eforall\egamma\enot\ephi$}
      \have[ ] {2} {$\enot\eexists\egamma\ephi$} \by{CQ}{1}
    \end{fitchproof}
  }{
    \begin{fitchproof}
      \have[n] {1} {$\enot\eexists\egamma\ephi$}
      \have[ ] {2} {$\eforall\egamma\enot\ephi$} \by{CQ}{1}
    \end{fitchproof}
  }
  \\
  and one pair for handling a negated universal quantifier: \\
  \twoitems{
    \begin{fitchproof}
      \have[n] {1} {$\eexists\egamma\enot\ephi$}
      \have[ ] {2} {$\enot\eforall\egamma\ephi$} \by{CQ}{1}
    \end{fitchproof}
  }{
    \begin{fitchproof}
      \have[n] {1} {$\enot\eforall\egamma\ephi$}
      \have[ ] {2} {$\eexists\egamma\enot\ephi$} \by{CQ}{1}
    \end{fitchproof}
  }
\end{fitchrule}

%\subsection{Rules for identity} \cite[][Chapter 39, pp. 326-9]{Mag25}

Finally, we will need some rules for handling the identity predicate ($=$). As
usual, these will take to form of introduction and elimination rules for $=$.

\begin{fitchrule}[Identity Introduction ($=$I)]
  \label{def:identity-introduction-rule}
  \leavevmode
  \vspace{-0.9\baselineskip}
  \begin{fitchproof}
    \have[ ] {1} {$\ealpha=\ealpha$} \by{$=$I}{}
  \end{fitchproof}
\end{fitchrule}

And, as with existential introduction, we can perform \emph{identity
elimination} on any one or more instances of the relevant term.

\begin{fitchrule}[Leibniz's Law ($=$E)]
  \label{def:identity-elimination-rule}
  Since identity is a symmetric relation, we define a pair of rules for the
  elimination of non-trivial identity sentences like $\ealpha=\edelta$. \\
  \leavevmode
  \vspace{-0.9\baselineskip}
  \twoitems{
    \begin{fitchproof}
      \have[m] {1} {$\ealpha=\edelta$}
      \have[n] {2} {$\ephi(\dots\ealpha\dots\ealpha\dots)$}    
      \have[ ] {3} {$\ephi(\dots\edelta\dots\ealpha\dots)$} \by{$=$E}{1,2}
    \end{fitchproof}
  }{
    \begin{fitchproof}
      \have[m] {1} {$\ealpha=\edelta$}
      \have[n] {2} {$\ephi(\dots\edelta\dots\edelta\dots)$}    
      \have[ ] {3} {$\ephi(\dots\ealpha\dots\edelta\dots)$} \by{$=$E}{1,2}
    \end{fitchproof}
  }
\end{fitchrule}

%\subsection{Derived rules **} \cite[][Chapter 40, pp. 330-1]{Mag25}

% TODO: Proofs of the CQ rules ??

%\subsection{Proofs and semantics **} \cite[][Chapter 41, pp. 332-5]{Mag25}

Without getting too deep into the metalogic, we can make some basic claims about
first-order logic. The most important observation is that FOL is both
\emph{sound} and \emph{complete}.

\begin{definition}[Soundness and Completeness of First-Order Logic]
  $ephi_1,\dots,\ephi_n\efollows\epsi$ iff
  $ephi_1,\dots,\ephi_n\eentails\epsi$  
\end{definition}

Corollary to this observation is the practical equivalence of several analogous
proof-theoretic and model-theoretic concepts.

\begin{definition}[Properties of First-Order Logic]
  Since first-order logic is sound and complete, the following pairs of concepts
  are analogues.\\
  \vspace{0.2cm}
  \begin{center}
  \begin{tabular}{c c c}
    \textbf{model-theoretic} & \textbf{iff} & \textbf{proof-theoretic} \\
    % \hline
    valid               & = & provable           \\
    validity            & = & theorem            \\
    equivalent          & = & interderivable     \\
    jointly satisfiable & = & jointly equivalent \\
  \end{tabular}
  \end{center}
\end{definition}

Strategies for proving these concepts in FOL metalogic can be summarized in the
usual table.

\begin{center}
  \begin{tabular}{c c c}
    & \textbf{yes} & \textbf{no} \\
    \hline
    &&\\    
    \emph{valid (sentence)?}    &
    prove $\efollows \ephi$ &
    find \ephi counter-model \\
    &&\\
    \emph{contradiction?}       &
    prove $\ephi\efollows\ebot$ &
    find $\enot\ephi$ counter-model \\
    &&\\    
    \emph{equivalent?}          &
    \makecell{prove $\ephi\efollows\epsi$ and\\$\epsi\efollows\ephi$} &
    \makecell{find model where \\\ephi and \epsi differ} \\
    &&\\    
    \emph{jointly satisfiable?} &
    \makecell{model where $\ephi_i$\\are all true} &
    prove $\ephi_1\dots\ephi_n\efollows\ebot$ \\
    &&\\    
    \emph{valid (argument)?}    &
    prove $\ephi_1\dots\ephi_n\efollows\epsi$ &
    \makecell{find model with $\ephi_1\dots\ephi_n$\\all true but \epsi false}\\
  \end{tabular}
\end{center}
